\section{Ease of Use}

\subsection{Maintaining the Integrity of the Specifications}



\section{Prepare Your Paper Before Styling}


\subsection{Abbreviations and Acronyms}



\section{Metodologia}

\subsection{Sistema Fuzzy Mamdani para Controle de Ventilação}

O sistema Mamdani foi projetado para decidir a intensidade da ventilação em ambientes fechados, considerando três variáveis de entrada: temperatura interna ($15^\circ$C a $35^\circ$C), umidade relativa (20\% a 100\%) e número de pessoas (0 a 20). Cada variável foi particionada em três conjuntos fuzzy com funções de pertinência triangulares~\cite{ross2010fuzzy}:

\begin{itemize}
    \item \textbf{Temperatura:} baixa, média, alta
    \item \textbf{Umidade:} baixa, média, alta
    \item \textbf{Pessoas:} poucas, moderado, muitas
\end{itemize}

A variável de saída, intensidade da ventilação, foi definida qualitativamente: fraca (0 a 4), moderada (2 a 8) e forte (6 a 10), também com funções triangulares.

A base de regras fuzzy foi construída com base em conhecimento comum sobre conforto térmico, utilizando operadores \textbf{AND} (mínimo), \textbf{OR} (máximo) e diferentes técnicas de agregação~\cite{ross2010fuzzy}. As principais regras são:

\begin{itemize}
    \item Se temperatura \textbf{alta} OU umidade \textbf{alta} OU pessoas \textbf{muitas} então ventilação \textbf{forte}
    \item Se temperatura \textbf{média} E umidade \textbf{média} E pessoas \textbf{moderado} então ventilação \textbf{moderada}
    \item Se temperatura \textbf{baixa} E umidade \textbf{baixa} E pessoas \textbf{poucas} então ventilação \textbf{fraca}
    \item (demais regras conforme detalhado no notebook)
\end{itemize}

A inferência fuzzy segue o modelo Mamdani~\cite{mamdani1975experiment}: fuzzificação das entradas, avaliação das regras, agregação das saídas e defuzzificação. Foram implementadas três técnicas de defuzzificação: centroide, média do máximo e bissetriz~\cite{ross2010fuzzy}.

\subsection{Sistema Fuzzy Takagi-Sugeno para Aproximação de Função Não Linear}

O sistema Takagi-Sugeno foi utilizado para aproximar a função $f(x) = e^{-x/5} \cdot \sin(3x) + 0.5 \cdot \sin(x)$ no intervalo $[0, 10]$. O sistema foi testado nas versões de ordem zero (consequente constante) e primeira ordem (consequente linear)~\cite{takagi1985fuzzy,sugeno1985industrial}.

Foram utilizadas três funções de pertinência gaussianas para a variável de entrada $x$, centradas em 2, 5 e 8, todas com $\sigma = 2$. As regras são do tipo:

\begin{itemize}
    \item Se $x$ é Baixo então $y = c_1$ (ordem zero) ou $y = a_1 x + b_1$ (primeira ordem)
    \item Se $x$ é Médio então $y = c_2$ ou $y = a_2 x + b_2$
    \item Se $x$ é Alto então $y = c_3$ ou $y = a_3 x + b_3$
\end{itemize}

A saída do sistema é a média ponderada das saídas das regras, ponderadas pelos graus de pertinência. Foram testadas funções de pertinência alternativas (triangular, trapezoidal) e operadores de agregação (produto, mínimo)~\cite{ross2010fuzzy}. Os consequentes foram otimizados usando gradiente descendente para minimizar o erro quadrático médio (MSE)~\cite{haykin2008neural}.

\subsection{Avaliação e Métricas}

Para ambos os sistemas, foram realizadas simulações com diferentes cenários de entrada e combinações de operadores. O desempenho do sistema Takagi-Sugeno foi avaliado usando as métricas MSE (Mean Square Error) e RMSE (Root Mean Square Error). Para o sistema Mamdani, foram analisadas as diferenças qualitativas e quantitativas entre as técnicas de defuzzificação e operadores lógicos.
